\documentclass[11pt]{article}
\usepackage{amsmath,amssymb,amsthm,hyperref,geometry}\geometry{margin=1in}
\newtheorem{theorem}{Theorem}
\title{Directed state extension with coherent pure actions}
\date{15 November 2026}\begin{document}\maketitle
\paragraph{Status.} Independently reviewed conditional integration theorem. It replaces the earlier symmetric-monoidal-source assumption by explicit pure-action and directed-parallel data. It does not infer these from an arbitrary relative pseudomonad. No strictification or free-feedback theorem is claimed.

\section{Independent process structure}
Let $C$ be a symmetric monoidal category, $\mathcal H$ a bicategory on the same objects, and $J:C\to\mathcal H$ the pure interface. Assume coherent left and right $C$-actions on $\mathcal H$, including their mixed associativity comparisons. The arrows $J(C)$ have specified centrality isomorphisms, natural in native two-cells and coherent under pure composition and tensor. These are actual Freyd/premonoidal action data, not an assertion that an unspecified centre is automatically monoidal.

Assume further parallel functors on process hom categories, a normal directed sequential compositor
\[
 \chi:(f\parallel g);(f'\parallel g')\Rightarrow(f;f')\parallel(g;g'),
\]
and coherent associativity, unit and symmetry comparisons for parallel. These comparisons satisfy the lax compositor equations and are compatible with $\chi$. Restrictions of parallel to pure/identity arguments are identified with the preceding actions by coherent invertible comparisons. In particular every structural move involving a pure state map and one effect uses those specified isomorphisms. No inverse of $\chi$ on two general effectful strands is supplied.

This is independent structure on the existing process bicategory. A relative parallel-composite theorem can construct part of it, but the identifications with the pure actions must still be verified; a bare lax parallel functor with unrelated actions is insufficient. Let $R:S\to C$ be a strong symmetric monoidal functor from a small state category.

\paragraph{A sufficient source of the pure actions.}
One need not supply unrelated centrality data. In a strict two-category presentation, assume $J$ is coherently parallel-strong, its pure composition constraints agree with $\chi$, and $\chi$ is invertible whenever one entire parallel strand consists of two pure arrows. Parallel associativity, unit and symmetry comparisons must respect $\chi$. Then holding an object is a pseudofunctor: its compositor has an entire identity strand. For an effect $f:A\to B$ and pure $h:C\to D$, the comparisons
\[
 (f\parallel1_C);(1_B\parallel Jh)\longrightarrow f\parallel Jh,
 \qquad
 (1_A\parallel Jh);(f\parallel1_D)\longrightarrow f\parallel Jh
\]
are invertible because their second strand is entirely pure. Their composite through the common target derives the centrality isomorphism. Naturality and compatibility with composition follow from the lax compositor equations by cancelling only these invertible pure-strand instances; the parallel structural compatibility gives the mixed action constraints. This does not invert the opposite-identity crossover for two arbitrary effects. The hypothesis is a restricted source presentation, not a blanket strictification of all bicategories.

\section{The construction and sequential proof obligation}
Use the directed balance hom categories of
\[
 F_{A,B}(s,t)=\mathcal H(JRs\otimes A,JRt\otimes B),
\]
where the displayed state whiskerings are the actual pure actions. For $[s,f]:A\to B$ and $[t,g]:B\to D$, keep both states, apply $f$ holding $Rt$, then $g$ holding $Rs$, and restore state order $s,t$. Thus chronological effect order is fixed.

On a slide in the first factor use the slide for $h\otimes1_t$; on a slide in the second use $1_s\otimes k$. The boundary comparisons are pure centrality and action isomorphisms. Native cells are whiskered by the action pseudofunctors and composed horizontally. The proposed associator uses the state associator slide and action coherence; it compares the same sequence $f,g,l$, not a permutation of effects.

\begin{theorem}[Action-based state extension]
Under the independent coherent data above, the construction gives a bicategory of directed state processes. Its pure-state unit inclusion is locally faithful. It inherits a directed parallel compositor, without inverting general effectful interchange.
\end{theorem}
\begin{proof}
The source/target actions of state maps define the pseudofunctor $F$ using their actual coherence. In the composition presentation, native/slide naturality is action naturality together with the chosen centrality of the pure state map. Two slides in one factor combine by the directed balance relation; a slide in each factor combines to the same state map $h\otimes k$. The remaining boundary maps agree by the mixed action and pure centrality equations. These checks establish descent on all product-category generators. They do not cancel a noninvertible slide.

For associativity, expand both three-process composites into held-state actions and pure wire maps. Apply the specified left/right and mixed action associators and centrality coherence to obtain the same chronological expression $f$ then $g$ then $l$. No step moves one effect past another. The four-process pentagon is the action associativity pentagon and mixed-action coherence, with the same pure-state permutation on both routes. Unitor equations use the action unit laws. Naturality on slides follows by normalizing only those boundary transformations induced from the same heterotyped body. This replaces, rather than invokes outside its scope, the earlier symmetric-monoidal-bicategory coherence theorem.

The unit inclusion uses state $I$ and the action unit isomorphisms. This inclusion is faithful, not generally full: completion can add comparisons even between unit-state processes. The original cells are preserved, not asserted to constitute every cell between their images. Its local faithfulness still follows from the intertwiner detector: a native cell is sent to the identical cell with unit boundaries, and the detector uses only the bicategory and state-map pseudofunctor.

For parallel, retain state $s\otimes t$, apply the source parallel body and reorder by pure structural maps. The inherited sequential compositor first reorders the four retained states, applies $\chi$ once to the two effectful time layers, then uses the pure-action identifications at its boundaries. Its direction is unchanged. Native and slide naturality follow from naturality of $\chi$ and its compatibility with the pure restrictions. The inherited compositor associativity/unit equations are the source lax equations with identical pure reorderings; the assumed compatibility of parallel structural comparisons supplies the remaining parallel associativity/symmetry coherence. Only pure boundary comparisons and state isomorphism slides are inverted.
\end{proof}

\section{State feedback, not an instantaneous trace}
Define $B_u[s,f]=[s\otimes u,f]$ for a process with interface $Ru\otimes A\to Ru\otimes B$, inserting the pure action associators. A native arrow is preserved and a slide $h$ becomes $h\otimes1_u$. This is a functor by the balance relations. An outer state map $u\to v$ induces the directed slide $1_s\otimes h$ on a heterotyped body.

\paragraph{Typed feedback comparisons.}
For $a:A'\to A$, $f:Ru\otimes A\to Ru\otimes B$, $b:B\to B'$ in the completion, use pure actions to form the displayed held-state composites. The comparisons are
\begin{align*}
 B_u((1_{Ru}\parallel a);f;(1_{Ru}\parallel b))
   &\cong a;B_u(f);b, &\text{tightening},\\
 B_I(f)&\cong f, &\text{vanishing},\\
 B_v(B_u(f))&\cong B_{u\otimes v}(f), &\text{joining},\\
 B_u(f)\parallel g&\cong B_u(\operatorname{reorder}(f\parallel g)),
   &\text{strength}.
\end{align*}
For joining, $f$ has interface $Ru\otimes Rv\otimes A\to Ru\otimes Rv\otimes B$. Reorder uses only pure associators and symmetries. If the internal state grades of $a,f,b$ are $r,s,t$, tightening changes the retained word $(r,s,t,u)$ to $(r,s,u,t)$. Strength changes $(s,u,t)$ to $(s,t,u)$. Joining reassociates $(s,u,v)$; vanishing removes $I$. The corresponding state isomorphism slide and pure-action boundary comparison define each cell.

\paragraph{Directed outer sliding.}
For $h:u\to v$ and a heterotyped $f:Rv\otimes A\to Ru\otimes B$, the comparison is
\[
 B_u(f\circ(JRh\otimes1_A))\longrightarrow
 B_v((JRh\otimes1_B)\circ f).
\]
It is the slide for $1_s\otimes h$ at the actual retained internal state $s$. Its naturality, identity and composition equations are balance equations, not inverse-slide equations.

\paragraph{Independent closure assembly and the exact interchanger square.}
Define $\operatorname{Seq}_{u,w}(f,f')$ to be the body on retained outer state $u\otimes w$ that applies $f$ holding $w$, then $f'$ holding $u$. Define $\operatorname{Par}_{u,v}(f,g)$ by parallel followed by the pure reordering that places outer states $u,v$ first. The canonical comparison maps are
\begin{align*}
 \Gamma_{\mathrm{seq}}:(B_u f);(B_w f')
 &\cong B_{u\otimes w}\operatorname{Seq}_{u,w}(f,f'),\\
 \Gamma_{\mathrm{par}}:(B_u f)\parallel(B_v g)
 &\cong B_{u\otimes v}\operatorname{Par}_{u,v}(f,g).
\end{align*}
If the internal states are $s,t$, these maps move $(s,u,t,w)$ to $(s,t,u,w)$, respectively $(s,u,t,v)$ to $(s,t,u,v)$, with their actual pure-action boundary cells. They do not merge state factors.

Take $f:Ru A\to Ru B$, $f':Rw B\to Rw E$, $g:Rv C\to Rv D$, $g':Rz D\to Rz F$ in the completion. Abbreviate
\[
 L=\operatorname{Seq}_{uv,wz}(\operatorname{Par}_{u,v}(f,g),
                                  \operatorname{Par}_{w,z}(f',g')),
\quad
 K=\operatorname{Par}_{uw,vz}(\operatorname{Seq}_{u,w}(f,f'),
                                  \operatorname{Seq}_{v,z}(g,g')).
\]
Products of state letters denote fixed bracketings of tensor. The claimed square is
\[
\begin{array}{ccc}
 ((B_uf)\parallel(B_vg));((B_wf')\parallel(B_zg'))
 &\xRightarrow{\chi_{\mathcal E}}&
 ((B_uf);(B_wf'))\parallel((B_vg);(B_zg'))\\
 \downarrow\scriptstyle{\Gamma_L}&&\downarrow\scriptstyle{\Gamma_R}\\
 B_{uv wz}(L)&\xRightarrow{\Xi}&B_{uw vz}(K).
\end{array}                                                   \tag{FC}
\]
Here $\Gamma_L$ first uses the two parallel assembly maps then sequential assembly; $\Gamma_R$ first uses the two sequential assembly maps then parallel assembly. To define $\Xi$, use the pure outer permutation $(u,v,w,z)\to(u,w,v,z)$ to transport the closure. On that outer order let $F$ be $f$ holding $w$, $F'$ be $f'$ holding $u$, $G$ be $g$ holding $z$, and $G'$ be $g'$ holding $v$. Pure-action identifications identify the transported $L$ with $(F\parallel G);(F'\parallel G')$. Apply the forward inherited compositor to obtain $(F;F')\parallel(G;G')$, identify this with $K$, and close under $B_{uwvz}$. These prescriptions define the lower arrow, including its state permutation, without an inverse of general $\chi$.

\begin{theorem}[Coherence of the stated state-feedback laws]
The four displayed isomorphisms and directed outer sliding are natural on the retained native hom categories. They satisfy tightening composition and unit laws, joining associativity and unit laws, strength associativity and unit laws, compatibility of strength with joining, and outer-slide compatibility with tightening, joining and strength. Moreover the explicitly defined square (FC) commutes, compatibly with the compositor associativity and unit equations.
\end{theorem}
\begin{proof}
On an internal cell, naturality is action/parallel naturality. On an inner state slide, both routes combine to the same tensor of the inner map with the identity outer state map; use balance naturality only on the corresponding heterotyped body. For an outer slide, the same argument combines to the tensor of the inner and outer maps. This verifies naturality without cancellation of a directed slide.

For tightening composition, both composites retain the same state list and execute the additional pre-processes, the middle process, and the post-processes in identical chronological order. The two native boundary pastings agree by pure-action associativity and centrality. Tightening units use the action unit laws. Joining associativity and units reduce to the state monoidal pentagon and triangles. Strength associativity and units are the parallel associator/unit equations together with the same pure permutations. Its joining compatibility permutes $(s,u,v,t)$ to $(s,t,u,v)$ in the two ways; symmetric pure coherence and the compatibility of parallel structure with the actions identify those maps. Outer-slide compatibility replaces the appropriate identity factors in these state permutations by $h$; naturality of the pure structural maps and balance composition prove the equations.

For (FC), expand the specified assembly maps. First use their state isomorphism slides to put all independently retained state factors in the order prescribed by $(u,v,w,z)\to(u,w,v,z)$. Both routes then apply the SAME source $\chi$ to the same two effectful layers. The remaining comparisons are pure-action identifications, so the source compatibility of $\chi$ with those identifications proves the square. For iterated layers use the source lax associativity equation, not an inverse of $\chi$. For iterated parallel factors use its compatibility with the parallel associator. This proves preservation of the directed compositor with its associativity and unit equations.
\end{proof}

\paragraph{The state-copy boundary.}
The preceding equations do not identify $B_u(f);B_u(g)$ with $B_u(f;g)$. The former retains two copies of $u$ and the latter one reused state. Sequential assembly of independently closed processes uses two distinct state factors and their held-state bodies. Any equation merging those factors would need extra diagonal/sharing data and a separate proof. The compositor compatibility above is only for the correctly typed independent factors.

As a separate sanity check, take the locally discrete process bicategory of sets and functions, finite-set states, and $u=\{0,1\}$. Let $f:u\times1\to u\times1$ flip the state, and $g:u\times1\to u\times u$ return the current state as output. The separate closure composite has retained state $u\times u$, while closure of the sequential body has retained state $u$. They cannot be isomorphic as objects of the directed completion's hom category: the intertwiner detector would send any such isomorphism to mutually inverse state maps, contradicting cardinalities four and two. If both state copies are initialized to zero in an additional interpretation, separate closure outputs zero while shared closure outputs one. Initialization is not assumed by the abstract construction. This test supports the stated boundary rather than defining the general theorem.

This states a specific structured-state feedback doctrine, not an instantaneous trace. A free universal property still requires its precise target morphisms and a separate evaluation theorem. No yanking law, arbitrary shared-interface admission, or completed research contract is asserted.

\paragraph{Literature.} Paquet--Saville, \emph{Premonoidal and Freyd bicategories}, \url{https://arxiv.org/abs/2312.14964}, relates Freyd structure to compatible monoidal actions. Their \emph{Effectful semantics in bicategories}, \url{https://philipsaville.co.uk/concurrent-pseudomonads.pdf}, Proposition 5.10 and Proposition 6.4, supplies the action and directed-parallel architectural precedents. These results do not automatically prove this directed-state extension or the broader feedback universal property. No novelty claim is made.
\end{document}
